% ═══════════════════════════════════════════════════════════════
% MT-06b-ML: Minitest MUSTERLÖSUNG - ir + estar
% Spanisch Klasse 7 | Sequenz 6: Mi ciudad
% ═══════════════════════════════════════════════════════════════

\documentclass[11pt, a4paper]{article}

\usepackage[utf8]{inputenc}
\usepackage[T1]{fontenc}
\usepackage[ngerman]{babel}
\usepackage{helvet}
\renewcommand{\familydefault}{\sfdefault}

\usepackage[margin=2cm]{geometry}
\usepackage{enumitem}
\usepackage{tabularx}
\usepackage{booktabs}

\usepackage{xcolor}
\usepackage{tcolorbox}
\tcbuselibrary{skins}

\usepackage{fancyhdr}
\usepackage{lastpage}
\usepackage{needspace}

% FARBEN
\definecolor{spanisch-color}{HTML}{c41e3a}
\definecolor{grau-color}{HTML}{888888}
\definecolor{hellgrau-color}{HTML}{f5f5f5}
\definecolor{loesung-color}{HTML}{c62828}

\colorlet{spanisch}{spanisch-color}
\colorlet{grau}{grau-color}
\colorlet{hellgrau}{hellgrau-color}
\colorlet{loesung}{loesung-color}

\newcommand{\fachfarbe}{spanisch}

\newcommand{\thema}{Minitest: ir + estar}
\newcommand{\fach}{Spanisch}
\newcommand{\klasse}{7}
\newcommand{\maxpunkte}{28}
\newcommand{\zeit}{15 Minuten}
\newcommand{\datum}{\today}

\pagestyle{fancy}
\fancyhf{}
\renewcommand{\headrulewidth}{0.4pt}
\renewcommand{\footrulewidth}{0.4pt}

\lhead{\fach{} -- Klasse \klasse}
\chead{\textcolor{loesung}{\textbf{MT-06b MUSTERLÖSUNG}}}
\rhead{\datum}

\lfoot{\zeit}
\cfoot{}
\rfoot{Seite \thepage\ von \pageref{LastPage}}

\newcommand{\punktebox}[1]{\hfill\fbox{\small \_\_\_/#1}}
\newcommand{\luecke}[1]{\rule{#1}{0.4pt}}
\newcommand{\lsg}[1]{\textcolor{loesung}{\textbf{#1}}}

\newtcolorbox{infobox}{
    colback=hellgrau,
    colframe=\fachfarbe,
    boxrule=1pt,
    arc=0pt,
    left=8pt, right=8pt, top=8pt, bottom=8pt
}

\newenvironment{aufgabe}[1]{%
    \needspace{5\baselineskip}%
    \nopagebreak[4]%
    \vspace{10pt}
    \noindent\textbf{\textcolor{\fachfarbe}{Aufgabe #1}}
    \nopagebreak[4]%
    \vspace{5pt}
    \nopagebreak[4]%
    \begin{list}{}{\leftmargin=0pt}
    \item[]
}{%
    \end{list}
}

\newcommand{\es}[1]{\textcolor{\fachfarbe}{\textbf{#1}}}

\begin{document}

% ═══════════════════════════════════════════════════════════════
% KOPFBEREICH
% ═══════════════════════════════════════════════════════════════

\begin{center}
    {\Large\bfseries\textcolor{\fachfarbe}{\thema}}\\[0.3em]
    {\large\textcolor{loesung}{\textbf{--- MUSTERLÖSUNG ---}}}
\end{center}

\vspace{5pt}

\noindent
\begin{tabularx}{\textwidth}{|X|X|X|}
\hline
\textbf{Name:} \lsg{Musterschüler/in} &
\textbf{Klasse:} \klasse &
\textbf{Datum:} \datum \\
\hline
\end{tabularx}

\vspace{10pt}

\begin{infobox}
\textbf{Zeit:} \zeit{} | \textbf{Punkte:} \maxpunkte{} BE | \textbf{Hilfsmittel:} keine
\end{infobox}

\vspace{15pt}

% ═══════════════════════════════════════════════════════════════
% AUFGABE 1: ir konjugieren
% ═══════════════════════════════════════════════════════════════

\begin{aufgabe}{1 -- ir konjugieren}
Ergänze die richtige Form von \es{ir}. \punktebox{6}
\end{aufgabe}

\noindent
\begin{tabular}{p{0.45\textwidth}p{0.45\textwidth}}
a) Yo \lsg{voy} al cine. & b) Tú \lsg{vas} a la plaza. \\[0.8em]
c) María \lsg{va} al parque. & d) Nosotros \lsg{vamos} a la biblioteca. \\[0.8em]
e) Ellos \lsg{van} al supermercado. & f) Vosotros \lsg{vais} al museo. \\
\end{tabular}

\vspace{15pt}

% ═══════════════════════════════════════════════════════════════
% AUFGABE 2: al vs a la
% ═══════════════════════════════════════════════════════════════

\begin{aufgabe}{2 -- al oder a la?}
Ergänze \es{al} oder \es{a la}. \punktebox{6}
\end{aufgabe}

\noindent
\begin{tabular}{p{0.45\textwidth}p{0.45\textwidth}}
a) Voy \lsg{al} supermercado. & b) Vas \lsg{a la} plaza. \\[0.8em]
c) Va \lsg{al} parque. & d) Vamos \lsg{a la} biblioteca. \\[0.8em]
e) Vais \lsg{a la} estación. & f) Van \lsg{al} hospital. \\
\end{tabular}

\vspace{0.5em}
\textcolor{loesung}{\textit{Regel: a + el = al (Kontraktion nur bei el!)}}

\vspace{15pt}

% ═══════════════════════════════════════════════════════════════
% AUFGABE 3: estar einsetzen
% ═══════════════════════════════════════════════════════════════

\begin{aufgabe}{3 -- estar einsetzen}
Setze die richtige Form von \es{estar} ein. \punktebox{8}
\end{aufgabe}

\noindent
a) El cine \lsg{está} en la plaza.

\vspace{0.8em}

\noindent
b) La biblioteca \lsg{está} cerca del parque.

\vspace{0.8em}

\noindent
c) ¿Dónde \lsg{estás} (tú)?

\vspace{0.8em}

\noindent
d) \lsg{Estoy} en casa. (yo)

\vspace{0.8em}

\noindent
e) Nosotros \lsg{estamos} en el parque.

\vspace{0.8em}

\noindent
f) Los museos \lsg{están} en el centro.

\vspace{0.8em}

\noindent
g) Vosotros \lsg{estáis} en la plaza.

\vspace{0.8em}

\noindent
h) María y Pedro \lsg{están} en la biblioteca.

\vspace{15pt}

% ═══════════════════════════════════════════════════════════════
% AUFGABE 4: ir vs estar
% ═══════════════════════════════════════════════════════════════

\begin{aufgabe}{4 -- ir oder estar?}
Kreise das richtige Verb ein. \punktebox{8}
\end{aufgabe}

\noindent
a) Yo \lsg{voy} / estoy al cine. (Ich gehe ins Kino) \textcolor{loesung}{$\rightarrow$ Bewegung!}

\vspace{0.8em}

\noindent
b) El cine va / \lsg{está} en la plaza. (Das Kino ist auf dem Platz) \textcolor{loesung}{$\rightarrow$ Lage!}

\vspace{0.8em}

\noindent
c) ¿Dónde va / \lsg{está} la farmacia? (Wo ist die Apotheke?) \textcolor{loesung}{$\rightarrow$ Lage!}

\vspace{0.8em}

\noindent
d) Nosotros \lsg{vamos} / estamos a la biblioteca. (Wir gehen...) \textcolor{loesung}{$\rightarrow$ Bewegung!}

\vspace{0.8em}

\noindent
e) Tú vas / \lsg{estás} en el parque. (Du bist im Park) \textcolor{loesung}{$\rightarrow$ Lage!}

\vspace{0.8em}

\noindent
f) Ellos \lsg{van} / están al hospital. (Sie gehen ins Krankenhaus) \textcolor{loesung}{$\rightarrow$ Bewegung!}

\vspace{0.8em}

\noindent
g) El museo va / \lsg{está} cerca. (Das Museum ist in der Nähe) \textcolor{loesung}{$\rightarrow$ Lage!}

\vspace{0.8em}

\noindent
h) ¿Adónde \lsg{vas} / estás (tú)? (Wohin gehst du?) \textcolor{loesung}{$\rightarrow$ Bewegung!}

% ═══════════════════════════════════════════════════════════════
% PUNKTEÜBERSICHT
% ═══════════════════════════════════════════════════════════════

\vspace{25pt}
\hrule
\vspace{10pt}

\noindent
\begin{tabularx}{\textwidth}{|X|c|c|c|c||c|}
\hline
\textbf{Aufgabe} & 1 & 2 & 3 & 4 & \textbf{Gesamt} \\
\hline
\textbf{max. Punkte} & 6 & 6 & 8 & 8 & \textbf{28} \\
\hline
\textbf{erreicht} & \lsg{6} & \lsg{6} & \lsg{8} & \lsg{8} & \lsg{28} \\
\hline
\end{tabularx}

\vspace{10pt}

\begin{flushright}
\textbf{Note:} \lsg{14 NP}
\end{flushright}

\vspace{1em}

\textbf{Notenspiegel (14-NP-Skala):}

\begin{center}
\footnotesize
\begin{tabular}{|c|c|c|c|c|c|c|c|c|c|c|c|c|c|c|}
\hline
\textbf{NP} & 14 & 13 & 12 & 11 & 10 & 9 & 8 & 7 & 6 & 5 & 4 & 3 & 2 & 1 \\
\hline
\textbf{ab BE} & 28 & 27 & 25 & 24 & 22 & 20 & 19 & 17 & 15 & 13 & 11 & 9 & 8 & 6 \\
\hline
\end{tabular}
\end{center}

\vspace{1em}

\textcolor{loesung}{\textbf{Merke:} ir = Bewegung (wohin?), estar = Lage (wo?)}

\end{document}
